\scanpage{12}\mspage{11}
\diagram{%
\[
\vcenter{\hbox{\begin{tikzcd}[ampersand replacement=\&,row sep=1.25em,column sep=.2em,
execute at end picture={\draw ([xshift=-.65em,yshift=1em]leq.north west) -- ([xshift=.65em,yshift=.1em]leq.north east) -- ([xshift=.75em,yshift=-.95em]ltlt.south east) -- ([xshift=-.9em,yshift=-.1em]ltlt.south west) -- cycle;}]
\& \& G\triangleleft F \arrow[d,Rightarrow] \& \& \\
G\preccurlyeq F \arrow[dr,Rightarrow] \& \& |[alias=leq]| G\leq F \arrow[dl,Rightarrow] \arrow[dr,Rightarrow] \& \& G\parallel F \arrow[dl,Rightarrow] \\
\& |[alias=ltlt]| G\ll F \& \& G\mathrel{\tikz[baseline=0pt,x=1em,y=1em]{\draw (-.38,.62)--(.38,.25)--(-.43,-.10); \draw (.34,.62)--(-.32,.25)--(.38,-.10);}} F \&
\end{tikzcd}}} \tag*{(1.18)}
\]}

Parmi les six relations, il y a quatre relations d'ordre $\triangleleft$, $\leq$, $\preccurlyeq$, $\ll$, et les deux autres relations $\parallel$ et $\mathrel{\tikz[baseline=0pt,x=1em,y=1em]{\draw (-.38,.62)--(.38,.25)--(-.43,-.10); \draw (.34,.62)--(-.32,.25)--(.38,-.10);}}$ sont symétriques.
% slanted
\margin{Non, $\triangleleft$ n'est pas rel. d'ordre, elle est transitive mais non réflexive, car $F\triangleleft F\Longleftrightarrow F\unsure{\in|\overset{\prime}{\in}}M$ Pour avoir une relation d'ordre il faudrait écrire \deleted{$\unsure{F\triangleleft|F}$} $F\trianglelefteq F\overset{\mathrm{déf}}{\Longleftrightarrow}F\triangleleft F$ \uuline{ou} $F=F$.}
La relation $\mathrel{\tikz[baseline=0pt,x=1em,y=1em]{\draw (-.38,.62)--(.38,.25)--(-.43,-.10); \draw (.34,.62)--(-.32,.25)--(.38,-.10);}}$ est \inserted{de plus} réflexive, mais $\parallel$ ne l'est que si $\unsure{\mathcal{F}|\mathscr{F}}$ est réduit à son seul \unsure{élémt|élément} $\varnothing_{\mathcal{F}}$, car
\[
 F\parallel F \quad\Longleftrightarrow\quad F=\varnothing_{\mathcal{F}}\qquad \text{figure vide.} \tag*{(1.19)}
\]

\uline{Remarques}\quad a) Les relations $\triangleleft$ et $\leq$ se déterminent mutuellement. Connaissant $\leq$, on définit $M\subset\mathcal{F}$ comme
\[
M=\left\{\begin{aligned}
&X\in\mathcal{F}\mid X\
\underset{\text{pour }\leq}{\text{$\mathrm{str}^{\mathrm{t}}$. irréductible}},\
\text{i.e.}\ X\leq\mathop{\mathrm{Sup}}\limits_i F_i\Rightarrow\\
&\hspace{12em}\exists i\ \text{avec }X\leq F_i
\end{aligned}\right\}
\tag*{(1.20)}
\]
puis $\triangleleft$ par (1.17).
\deleted{\unsure{Et|En}} D'autre part, connaissant $\triangleleft$ on définit $M$ par
\[
M=\{X\in\mathcal{F}\mid X\triangleleft X\} \tag*{(1.21)}
\]
puis, pour $F\in\mathcal{F}$, son \unsure{déployement|déploiement}
\[
\widetilde F=\{X\in M\mid X\triangleleft\unsure{\mathcal{F}|F}\} \tag*{(1.22)}
\]
et enfin
\[
F\leq G\quad\Longleftrightarrow\quad\widetilde F\subset\widetilde G \tag*{\unsure{(1,23)|(1.23)}}
\]
Donc, pour déterminer \deleted{la structure, il est \unsure{équivalent}} \inserted{les données \unsure{minima|minimes} qui} fixent la structure d'atelier, il est indifférent d'y
